\documentclass[10pt,a4paper]{article}
\usepackage{cvcommon}
\hypersetup{pdftitle={CV ATS Javier López Molinero - Español},pdfauthor={Javier López Molinero},pdfsubject={Currículum ATS}}

\begin{document}
\ATSSetup
\ATSHeader{Senior C++ / Quant Developer · DevOps}{Madrid, España}{Experiencia profesional en equipos y proyectos internacionales}

\section{Perfil profesional}
Ingeniero de Telecomunicaciones con más de 15 años de experiencia en desarrollo de software, integración de sistemas e infraestructura CI/CD. Especializado en C++, Python, Java, DevOps, sistemas de tiempo real, microservicios y Linux.

\section{Competencias principales}
\textbf{Lenguajes:} C++, C, Python, Java, ADA, PHP, TTCN-3.\\
\textbf{Plataformas:} Linux, Yocto Linux, Unix, Windows Server, macOS.\\
\textbf{DevOps:} CI/CD, Docker, Kubernetes, Jenkins, Gerrit, testing, Scrum.\\
\textbf{Especialización:} integración de sistemas, cloud, SaaS, IoT, firmware embebido, tiempo real, TCP/IP, OSI, SDN y NFV.\\
\textbf{Idiomas:} inglés con competencia profesional y experiencia en entornos internacionales.

\section{Experiencia profesional}
\ATSJob{Quant Developer SR and Integrator Engineer for BBVA}{ARENA Financial Tech}{Abr. 2026 - Actualidad}{Desarrollo de software y evolución de la infraestructura de integración continua para la Sala de Tesorería de BBVA, reforzando la mantenibilidad y fiabilidad del entorno de Quants.}
% METRICA: añadir reducción del tiempo de compilación, frecuencia de despliegue o incidencias evitadas.
\ATSJob{Ingeniero desarrollador e integrador para BBVA}{Optimissa-Alten}{Mar. 2022 - Mar. 2026}{Desarrollo y mantenimiento de componentes C++ y Python, automatización CI/CD y tareas DevOps para Quants Develop Design and Cloud en la Sala de Tesorería.}
\ATSJob{Ingeniero desarrollador e integrador para Telefónica}{Full On Net}{Nov. 2019 - Mar. 2022}{Diseño y mantenimiento del entorno CI/CD de Movistar IPTV con C++, Python, servidores, redes, migración de entornos y Yocto Linux.}
\ATSJob{Advanced Development Engineer}{Indra · Proyecto HERMES / HORUS}{Jul. 2018 - Nov. 2019}{Desarrollo del centro de control urbano con C++, Java y Python; microservicios, Docker, Kubernetes, automatización DevOps y scripting.}
\ATSJob{Ingeniero de Desarrollo}{Altran · Cliente Ericsson}{Oct. 2016 - Jul. 2018}{Desarrollo C++ y TTCN-3 e integración continua para el nodo HSS en un equipo multidisciplinar de I+D.}
\ATSJob{Ingeniero de Software}{Indra Sistemas / Novanotio}{Feb. 2012 - Oct. 2016}{Firmware para sistemas de control de tráfico vial y coordinación del desarrollo del software del radar Doppler.}
\ATSJob{Consultor Junior en Sistemas}{TB-Security}{Abr. 2011 - Ago. 2011}{Consultoría de sistemas y soporte técnico.}
\ATSJob{Gestión tecnológica y soporte Moodle}{CIFF - Banco Santander}{Nov. 2010 - Abr. 2011}{Administración y soporte de la plataforma tecnológica de formación.}
\ATSJob{Preventa, soporte, comercial y formación}{Funkwerk EC Iberia}{Mar. 2010 - May. 2010}{Soporte y homologación de soluciones de comunicaciones.}
\ATSJob{Servicios IT y telecomunicaciones}{Quantum 2003, S.L.}{Jun. 2004 - Oct. 2009}{Servicios de informática, redes y telecomunicaciones.}

\newpage
\section{Formación}
\textbf{Máster en Ingeniería de Telecomunicaciones}, UOC, 2014-2018.\\
\textbf{Ingeniería Técnica de Telecomunicaciones - Telemática}, Universidad de Alcalá, 2003-2009.\\
\textbf{Técnico Superior en Sistemas de Telecomunicaciones e Informáticos}, La Salle-Sagrado Corazón, 2001-2003.

\section{Formación complementaria}
Programación Java (250 h), 2012 · CCNA, ESFgroup, 2011 · Cisco CCNA 1 y 2, University of NEWI (Wales), 2007.

\section{Proyectos relevantes}
\begin{itemize}
  \item \textbf{BBVA - Sala de Tesorería:} desarrollo de módulos y mejora continua de infraestructura CI/CD para Quants.
  \item \textbf{Movistar IPTV:} entorno CI/CD y jaula de compilación para descodificadores con C++, Python y Yocto Linux.
  \item \textbf{HERMES / HORUS:} back-end del centro de control semafórico de Indra con C++, Java y servicios distribuidos.
  \item \textbf{Nodo HSS - Ericsson I+D:} microservicios, integración continua y Kubernetes.
  \item \textbf{Control de tráfico:} firmware para radares y coordinación del software del radar Doppler.
  \item \textbf{Radio online laisla.fm:} diseño, creación e implantación comercial de la plataforma y de un reproductor para iPhone con más de 1.200 descargas.
  \item \textbf{Proyecto fin de máster:} MQTT-S sobre IEEE 802.15.4e en OpenMote, UOC, 2018. Prueba de concepto IoT con OpenMote, Raspberry Pi, sensores, base de datos en la nube, Grafana y OpenWSN.
\end{itemize}

\section{Aficiones}
Programación y proyectos personales · domótica · karate · inteligencia artificial aplicada a la mejora de procesos y al aprendizaje.
\end{document}
